Table~\ref{tab:ratio} varies the false-negative to false-positive regret of the reward from 1:1 to 10:1 (E-DDQN configuration per fold fixed; ratio 5:1 reproduces the main E-DDQN run bit for bit, which also confirms the determinism of the pipeline). The greedy policy moves along the trade-off exactly as Eq.~\eqref{eq:margin} predicts: recall rises from \RatioArgOneRecall\,\% (1:1) to \RatioArgFiveRecall\,\% (5:1) and \RatioArgOneZeroRecall\,\% (10:1), while specificity falls from \RatioArgOneSpec\,\% to \RatioArgFiveSpec\,\% and \RatioArgOneZeroSpec\,\%. Once the threshold is chosen on validation data, however, all ratios end at statistically indistinguishable costs (\RatioTunedOneCost--\RatioTunedOneZeroCost) and similar ranking: the reward ratio sets the default operating point, but it does not create discrimination that a threshold cannot reproduce. The 5:1 ratio is therefore best understood as a declared safety preference whose operating point can equally be reached by cost-weighted training or threshold selection.

\begin{table*}[t]
\centering
\caption{Reward-ratio sensitivity of E-DDQN: operating point of the greedy (arg-max) policy, whose implied cut-off on $P(y=1\mid s)$ is $1/(1+R)$, and of the validation-tuned threshold (rule D1). Mean ($\pm$ SD) over five seeds.}
\label{tab:ratio}
\footnotesize
\begin{tabular}{lcccccccc}
\toprule
 & & \multicolumn{3}{c}{Greedy policy ($\arg\max_a Q$)} & \multicolumn{3}{c}{Validation-tuned threshold (D1)} & \\
FN:FP regret & implied cut-off & Recall & Spec. & Cost$_{5:1}$ & Recall & Spec. & Cost$_{5:1}$ & ROC-AUC \\
\midrule
1:1 & $1/2$ & 52.4$\pm$5.6 & 58.1$\pm$4.8 & 1.285 & 94.9 & 7.4 & 0.629$\pm$0.056 & 0.558 \\
2:1 & $1/3$ & 69.2$\pm$3.3 & 43.1$\pm$1.6 & 0.999 & 94.6 & 6.4 & 0.642$\pm$0.074 & 0.563 \\
5:1 & $1/6$ & 77.2$\pm$3.8 & 28.0$\pm$2.0 & 0.905 & 93.4 & 7.4 & 0.663$\pm$0.048 & 0.538 \\
10:1 & $1/11$ & 80.7$\pm$3.2 & 23.0$\pm$3.3 & 0.855 & 92.9 & 7.4 & 0.673$\pm$0.095 & 0.552 \\
\bottomrule
\end{tabular}

\end{table*}
